\vspace{-1mm}
\section{Proposed Method}
\label{sec:methodology}



\vspace{-1mm}
\begin{figure*}[h!]
\centering
    \includegraphics[width=\textwidth]{Figures/system_overview.pdf}
    \vspace{-2mm}
    \caption{
        Overview of \texttt{HELIOS}:
        The initial graph $G_{init}$ is early-fused to generate a graph $G_d$.
        Each node and edge of $G_d$ are embedded.
        (1) The edges of $G_d$ are retrieved using the query $q$, then integrated into a candidate bipartite subgraph $G_c$.
        (2) The most query-relevant nodes in $G_c$ are identified as seed nodes.
        % Nodes of $G_{init}$ that are relevant to both the seed node and the query are expanded into $G_c$, forming the expanded graph $G_l$.
        New nodes from $G_{init}$ are expanded from the seed nodes, forming an expanded graph $G_l$.
        (3) LLM performs aggregation over restored tables to identify new relevant table rows, and then eliminates irrelevant passages.
    }
    \vspace{-5mm}
\label{fig:system_overview}
\end{figure*}
We propose \texttt{HELIOS}, a novel graph-based retrieval framework that combines the strengths of early fusion, late fusion, and LLM reasoning. 
As shown in Figure~\ref{fig:system_overview}, it operates in three stages: 
(i) \textbf{Edge-based Bipartite Subgraph Retrieval} retrieves edges from a bipartite data graph constructed via early fusion and integrates them into a single bipartite subgraph. 
(ii) \textbf{Query-relevant Node Expansion} enhances the retrieved subgraph by incorporating additional nodes through further retrieval. 
(iii) \textbf{Star-based LLM Refinement} refines the expanded graph through aggregation and multi-hop reasoning using the LLM.




% We propose \texttt{HELIOS}, a novel graph-based retrieval to leverage the advantages of both early and late fusion. 
% It operates in three main stages, as illustrated in Figure~\ref{fig:system_overview}:

% \textbf{Edge-based Bipartite Subgraph Retrieval.}
% It retrieves the edges of a bipartite data graph generated by early fusion, then integrates the retrieved edges into a single bipartite subgraph.



% \textbf{Query-relevant Node Expansion.}  
% It reinforces the retrieved bipartite subgraph with new nodes found through additional retrieval.  
% % In this step, we identify seed nodes, the base nodes for additional retrieval.  
% % Seed nodes are combined with the query to generate expanded queries, which retrieve relevant candidate nodes.  
% % The most promising candidates are then expanded into the subgraph.



% \textbf{Star-based LLM Refinement.}
% % It further refines the reinforced graph by perfoming aggregation and multi-hop reasoning via LLM's advanced reasoning.
% It refines the expanded graph through aggregation and multi-hop reasoning using the LLM.
% It first performs a conditional aggregation operation; restoring the original tables of the table segments within the expanded graph, then using LLMs to perform actual aggregation operations.
% First, when aggregation is determined necessary, it restores original tables from table segments in the expanded graph and applies aggregation via the LLM.
% The graph is updated with table segments identified as aggregation outputs, along with their related passages.
% Second, it removes the edges that are irrelevant in answering the query for each star graph within the graph.
% The edges verified as irrelevant are excluded.

















\vspace{-1mm}
\subsection{Edge-based Bipartite Subgraph Retrieval}
\label{sec:candidate_bipartite_graph_generation}
\vspace{-1mm}


% 1. 세부 스텝 두 가지에 대한 소개
\texttt{HELIOS} initiates its process with the retrieval of a bipartite subgraph through two key steps: \textit{early fusion} and \textit{edge retrieval}.


% 2. Early fusion step에 대한 설명
% 3. Edge & node embedding generation
\textbf{(1) Early fusion}: 
        % Query가 들어오기 전
This step is performed offline, before a query is given.
A bipartite data graph $G_d$ is constructed by generating edges from $G_{init}$, which initially has no edges. % by linking table segments and passages via entity linking.
% Each edge in this graph represents a meaningful connection between the table segment and the passage. 
% An embedding is also computed for each edge in this step.
% The step starts by generating edges from the initial graph $G_{init}$ that has no edges.
Edge generation is composed of entity recognition and entity linking, following prior methods~\citep{CORE, COS}.
% The output graph of this process $G_d$ is a bipartite graph, since the edges are generated only between passage nodes and table segment nodes.
Edges are generated between passage nodes and table segment nodes, resulting the generated data graph $G_d$ to be a bipartite graph.
% The next step is to generate embeddings for each edge.
% Previous early fusion techniques create a fixed-dimensional single embedding for each star graph. 
% In contrast,
% We employ a late interaction model (\citealp{ColBERTv2}) that generates token-level embeddings at the edge level, enabling richer information representation while reducing information loss. 
% Through experiments, we demonstrate that the late interaction model is well-suited for retrieval tasks involving large units, such as edge-level combinations of two documents \ref{}.
$G_d$ is then indexed in an edge-wise manner.
% The generated edges are are embedded into a sequence of vectors.
Each edge $e = (S, P)$ is first linearized into a token sequence $x$.
\vspace{-1mm}
\begin{equation}    
    x = [\ Linearize(\Gamma(S));\ \Gamma(P)\ ]
    \vspace{-1mm}
\end{equation}
Then, $x$ is embedded into a sequence of vectors.
% Mathematically, the encoding of both the query $q$ and the token sequence $x$ can be expressed as:
\vspace{-2mm}
\begin{equation}    
    \textbf{X} = f_{e}(x) \in \mathbb{R}^{l_x \times d}
    \vspace{-2mm}
\end{equation}
where $l_x$ represents the length of a linearized edge $x$.
% The function $f_e$ denotes our late-interaction edge encoder, which is fine-tuned following the training scheme of \texttt{ColBERTv2} (\citealp{ColBERTv2}).
{\editcolor We adopt the multi-vector encoder \texttt{ColBERTv2} (\citealp{ColBERTv2}) for $f_e$ to create fine-grained embeddings.
That is, HELIOS embeds edges, a larger unit compared to the node-level embeddings used in previous methods (\citealp{CORE, COS}). 
Since edges contain more tokens than nodes, fine-grained embeddings are essential to reduce information loss.}
% Detailed explanations of the fine-tuning process, including the construction of the training dataset, can be found in Appendix~\textsection~\ref{sec:edge_retriever_and_reranker}



\textbf{(2) Edge Retrieval}: 
% 4. Edge retrieval
% In this step, the set of most query-relevant edges is identified by leveraging the semantic similarity between the query and the edge embeddings.
        % Query가 들어옴.
When the query is given, we first identify the query-relevant edges by leveraging the semantic similarity between the query and edge embeddings.
% The edges are integrated to construct the candidate bipartite subgraph $G_c \subset G_d$.
% At online time, a bipartite candidate subgraph is generated by retrieving and merging the offline-embedded edges.
% This step is conducted through a three-stage process.
% An initial retrieval of edges is done using late interaction. % edge encoder %$f_{e}$ computes the similarity scores between the query $q$ and each edge $e$. 
% We use the same encoder $f_e$ to encode the query $q$.
\vspace{-2mm}
\begin{equation}
    \textbf{Q} = f_{e}(q) \in R^{l_q \times d}
    \vspace{-2mm}
\end{equation}
The similarity is then calculated between the query and each indexed edge.
\vspace{-2mm}
\begin{equation}
     sim(q, e;f_{e}) = \sum_{i=1}^{l_q} \max_{j\in [1, l_x]} \textbf{Q}_i\textbf{X}_j^T
     \vspace{-2mm}
\label{eqn:similarity_query_edge}
\end{equation}
We then select the top-$k_1$ edges that show the highest score $sim(q, e;f_{e})$, measuring query-edge alignment.
The selected edges are further passed to a fine-tuned all-to-all interaction reranker $g_{e}$, which performs a more detailed similarity evaluation.
This identifies the most contextually relevant edges, allowing us to identify the top-$k_2$ query-relevant edges $(k_2 < k_1)$.
% The details for the fine-tuning $g_e$ is also explained in Appendix~\textsection~\ref{sec:edge_retriever_and_reranker}.
The fine-tuning process for $f_e$ and $g_e$, including training dataset construction, is detailed in Appendix~\textsection~\ref{sec:edge_retriever_and_reranker}.


The resultant edges are integrated into the bipartite subgraph $G_c = (V_c, E_c, \Phi, \Gamma, \Lambda_c)$. 
Specifically, if there are duplicate nodes contained in the retrieved edges, those nodes are merged to form a single graph.
$G_c$ serves as the candidate bipartite subgraph, which becomes the foundation for further expansion and refinement.
We save the similarity score for each edge in the score mapping function $\Lambda_c$, which is later used to sort the edges based on their relevance to the query in Section~\ref{sec:graph_refinement_via_complex_reasoning}.
% $\Lambda_c$ will not be used until Section~\ref{sec:graph_refinement_via_complex_reasoning} where the final ranked list of edges will be generated.

% 5. Training the models
% The late interaction encoder $f_e$ , and the all-to-all interaction reranker $g_e$ was fine-tuned using contrastive loss.
% Detailed explanations of the fine-tuning process, including the construction of the training dataset, can be found in Appendix~\textsection~\ref{sec:edge_retriever_and_reranker}.


















\vspace{-1mm}
\subsection{Query-relevant Node Expansion}
\label{sec:candidate_bipartite_graph_augmentation}
\vspace{-1mm}



\begin{figure*}[h!]
\centering
    \includegraphics[width=0.85\textwidth]{Figures/Query-relevant_Node_Expansion.pdf}
    \caption{
        The overall procedure of query-relevant node expansion.
        The beam width $b$ is set as 2 in this example. 
        The purple-colored nodes indicate the selected seed nodes.
    }
    \vspace{-5mm}
\label{fig:query_driven_late_fusion}
\end{figure*}



% 1. Query-relevant node expansion
Query-relevant node expansion process aims to identify additional query-relevant edges, including the edges that have not been present within $G_d$. 
We perform the expansion process at the node level, which is the most fine-grained level. 
This is to address the issue that early fusion inevitably includes query-irrelevant nodes in the candidate subgraph.
% This on-the-fly expansion generates graph $G_l$, a graph that includes additional expanded nodes, from $G_c$.
% We perform the expansion process at the node level, which is the most fine-grained level. 
% This is to address the issue that early fusion inevitably includes query-irrelevant nodes in the candidate subgraph as the fused blocks are determined independent of the query.
% We choose node as our unit of retrieval in this process, ~~~.
We formalize the process as finding a set of edges that meet the following objective function:
% Formally, the process can be expressed as finding a set of edges that meet the following objective function:
\vspace{-3mm}
\begin{equation}
    \argmax_{(u,v)\in E^*\wedge u \in V_c} p(u,v|q) =  p(v|u,q)p(u|q)
    \vspace{-2mm}
\label{eqn:query_driven_late_fusion_objective}
\end{equation}
Here, $u$ represents a node in the candidate graph $G_c$, and $v$ represents a node adjacent to $u$ in the complete bipartite graph $G^*$. 
The complete bipartite graph $G^* = (\mathcal{C}_{init}, E^*, \Phi, \Gamma, \Lambda_{init})$ contains all possible edges between table segments and passages. 
% We employ a two-step beam search to identify expanded edges.
Naively solving this objective requires calculating the similarity score between the query and every possible edge in the complete graph, which is infeasible. 
To address this, we adopt a beam search approach and model it as a two-step process, as illustrated in Figure~\ref{fig:query_driven_late_fusion}.



% 2. Steps for query-relevant node expansion
% 3. Seed node selection
\textbf{(1) Seed node selection}: 
This step composes the first iteration of the beam search, corresponding to finding the set of nodes that show the highest $p(u|q)$.
% From the candidate bipartite subgraph $G_c$, we select a set of nodes that contain information most relevant to the query.
We calculate $p(u|q)$ for each $u \in V_c$ to identify the top-$b$ (i.e., beam width) nodes that contain the most relevant information to the query.
The probability $p(u|q)$ is determined by calculating the semantic similarity between the query and each node $u$ in $G_c$, then normalizing the scores using a softmax function. 
This similarity is computed through a fine-tuned all-to-all interaction-based node reranker $g_{n}$.
% \vspace{-2mm}
% \begin{equation}
%     p(u|q) \approx g_{n}([q ; \Gamma(u)])
%     \vspace{-2mm}
% \end{equation}
 

    % 4. Expanding node selection
\textbf{(2) Seed node expansion}: 
% For each seed node, we iterate through its neighbors in the complete bipartite graph $G^*$, calculating the similarities between each expanding node and the pair of the query and the seed node. 
% Among these, node pairs that exhibit the highest similarity with the query are returned as edges, further expanding the graph $G_c$ to $G_l$ by adding these edges.
This step composes the second iteration of the beam search, aiming to find the set of edges $(u, v)$ which show the highest $p(u, v|q)$.
It is done by computing $p(v|u,q)$ for each node $v$ connected to a seed node $u$ in the complete bipartite graph $G^*$.
This conditional probability is calculated using the expanded query retrieval technique (\citealp{expandedqueryretrieval}).
In this technique, the score function is expressed based on the expanded query as $sim([q;\Gamma(u)], v)$, and it is calculated by two late interaction models: 
$sim([q;\Gamma(u)], v; f_{P \rightarrow S})$
 for a table-segment-typed expanding node and a passage-typed seed node, and $sim([q;\Gamma(u)], v; f_{S \rightarrow P})$ for the opposite.
% {\color{red}
% \begin{equation}
%     sim_e([q;\Gamma(u)], v) = 
%     \begin{cases}
%           \sum_{i=1}^{l_q + l_{\Gamma(u)}} \max_{j \in [1, l_v]} \textbf{Q}_i \textbf{V}_j^T, & \text{if } i = j, \\
%           -y_{ij}, & \text{if } i \neq j \text{ and } i \sim j.
%     \end{cases}
% \end{equation}
% }
The calculated scores are normalized using a softmax function to compute $p(v|u,q)$.
Finally, $p(u, v|q)$ is calculated for the $(u, v)$ pairs following Equation~\ref{eqn:query_driven_late_fusion_objective}, and the top-$b$ edges with the highest scores are finally selected.
The $(u, v)$ pairs are added to $E_c$ along with the nodes $v$ to $V_c$, forming the expanded bipartite graph $G_l = (V_l, E_l, \Phi, \Gamma, \Lambda_l)$. 
The scores for each new edge are calculated using the reranker $g_e$ as used in \textsection~\ref{sec:candidate_bipartite_graph_generation}.
Detailed explanation for fine-tuning $g_n$, $f_{P\rightarrow S}$ and $f_{S\rightarrow P}$ can be found in Appendix~\textsection~\ref{sec:node_reranker} and \textsection~\ref{sec:expanded_query_retrievers}.

% The change in retrieval accuracy of \texttt{HELIOS} based on beam width $b$ is discussed in Appendix~\textsection~\ref{sec:exp_parameter_sensitivity}.
% The node reranker $g_{n}$ are fine-tuned for node selection.
% The late interaction encoders $f_{P\rightarrow S}$ and $f_{S\rightarrow P}$ used for expanded query retrieval are fine-tuned following the \texttt{ColBERTv2} training scheme (\citealp{ColBERTv2}).
% Detailed explanations for fine-tuning, including the construction of the training dataset, for all the modules can be found in Appendix~\textsection~\ref{sec:node_reranker} and Appendix~\textsection~\ref{sec:expanded_query_retrievers}.



















\vspace{-2mm}
\subsection{Star-based LLM Refinement}
\label{sec:graph_refinement_via_complex_reasoning}
\vspace{-1.5mm}

% Traditional semantic similarity is insufficient to correctly retrieve documents for queries that require logical inference, such as column-wise aggregation or multi-hop reasoning.
% To overcome this problem, we leverage the advanced reasoning capabilities of large language models (LLMs), which allow us to refine the retrieved result through logical inference.
% 목표: 그래프를 마지막으로 refine하는 것

% LLM의 추론 능력을 활용하여 semantic similarity만으로는 찾기 어려운 관련정보가 검색될 수 있도록 하는 것이다.
The main goal of this step is to retrieve query-relevant information which is challenging to find using semantic similarity alone, by leveraging the logical inference capabilities of LLMs.
% An overview of the step can be found in Appendix~\ref{sec:star_based_llm_refinement_supp}.
Selecting the optimal format or unit for presenting the graph $G_l$ to the LLM is non-trivial.
% It is non-trivial to choose the specific format or unit of providing the graph $G_l$ to the LLM. 
We explored two approaches: (i) providing the entire graph $G_l$ in a single prompt to return the relevant set of nodes and (ii) decomposing $G_l$ into star graphs, with each star graph generating its own set of relevant nodes.
The latter approach proved to be 12.4\% more effective, leading us to adopt star graphs as the unit for logical inference (\textsection~\ref{sec:exp_granularity_analysis}). 
An overview of the process and few-shot prompts is in Appendix~\ref{sec:star_based_llm_refinement_supp} and \textsection~\ref{sec:Prompts}.
The process consists of two phases: \emph{column-wise aggregation} and \emph{passage verification}.
% It is non-trivial to choose the specific format or unit of providing the graph $G_l$ to the LLM.
% We considered two approaches: including the entire graph $G_l$ in a single prompt to return the relevant set of nodes, and decomposing $G_l$ into star graphs, with each star graph generating its own set of relevant nodes. 
% Among these, using star graphs as the unit proved to be 12.4\% more effective, leading us to select this as our unit for logical inference (\textsection~\ref{sec:exp_granularity_analysis}).

% It is non-trivial to choose the specific format or unit of providing the graph $G_l$ to the LLM.
% We considered two approaches: including the entire graph $G_l$ in a single prompt to return the relevant set of nodes, and decomposing $G_l$ into star graphs, with each star graph generating its own set of relevant nodes. 
% Among these, using star graphs as the unit proved to be 12.4\% more effective, leading us to select this as our unit for logical inference (\textsection~\ref{sec:exp_granularity_analysis}).
% With this, the process occurs in two phases: \emph{column-wise aggregation} and \emph{passage verification}.

%%%%%%%%%%%%%%%%%%%%%%%%%%%%%%%%%%%% WORKING HERE %%%%%%%%%%%%%%%%%%%%%%%%%%%%%%%%%%%%


\textbf{(1) Column-wise Aggregation}: 
% The column-wise aggregation step restores tables from table segments, then identifies the candidate rows based on the query and adds them back to the graph.
This step aims to infer the correct result rows for table aggregation operations, as exemplified in Figure~\ref{fig:motivation}(c), making it possible that the corresponding rows exist in $G_l$. 
% In our early fused graph $G_d$, tables are divided into individual row-based table segments, making it hard for the previous-stage retrievers to perform proper aggregation.
% It becomes necessary to reconstruct the original tables and perform reasoning over these full tables.
Since not every query requires aggregation, we first prompt the LLM to determine whether the input query necessitates an aggregation operation. 
If the query is classified as an aggregation query, it first restores the original from each table segment using mapping function $\mathcal{M}$. 
The restored tables are then provided to the LLM in the format of star graph along with the query. 
The LLM performs the aggregation and returns the rows corresponding to the aggregation result.
The returned rows (i.e., table segments) are subsequently added to $G_l$ along with their associated passages to generate $G_a$.


\textbf{(2) Passage Verification}: 
The edges of $G_a$ may contain lots of hard negatives, as they comprise the edges retrieved from the edge retrieval, node expansion and column-wise aggregation step.
The passage verification step aims to identify the passages relevant to answering the query.
% The refined edge-scored graph is then decomposed into a ranked sequence of edges to produce the final retrieval output.
% It aims to leverage the LLM's logical inference capabilities to remove the query-irrelevant passages within $G_a$.
Similar to the column-wise aggregation step, we provide $G_a$ to the LLM in the form of star graphs, units that contain multi-hop relationships.
The LLM performs a binary verification to determine whether each edge is relevant to the query, without recalculating their scores. 
As a result, query-irrelevant edges are removed, yielding a refined edge-scored graph $G_q$. 
% Examples of prompts for these tasks are in Appendix~\textsection~\ref{sec:Prompts}.

We transform the graph $G_q$ into a sequence of edges, as our reader requires a serialized token sequence as input.
% The graph $G_q$ is then transformed into a ranked sequence of edges $\mathcal{E}_q$ by applying the score mapping $\Lambda_c$. 
We use the mapping function $\Lambda_c$ to sort each edge $e$ by its similarity to the query $\Lambda_c(e)$, and the top $k$ edges are returned.

